\chapter{Introdução}\label{ch:intro}
Este modelo de documento foi elaborado para atender às necessidades dos estudantes dos cursos de graduação do IFPE Campus Garanhuns.
Embora a biblioteca do IFPE forneça um modelo, considero sua utilização improdutiva para trabalhos acadêmicos de nível superior (graduação e pós-graduação). Não somente os aspectos diretamente relacionados à formatação, mas também em relação à estrutura do documento no tocante às seções, subseções, inserções de imagens ou tabelas com respectivas legendas e referências cruzadas são mais fáceis uma vez que se compreende o funcionamento do \LaTeX.

Também considero mais fácil o gerenciamento de referências bibliográficas e suas respectivas citações ao longo do texto, algo que no Word ou LibreOffice, por exemplo, deve ser feito com um programa externo (Zotero ou Mendeley) e seus respectivos \textit{plugins} para esses editores de texto.

Esse modelo obedece às normas ABNT eestá disponível gratuitamente no GitHub. Para aqueles que desejam realizar alguma modificação ou melhoria no modelo, sugiro que realize as mudanças diretamente no arquivo \textbf{definicoes.tex} a fim de manter a estrutura e organização do projeto. Veja o README.md no Github para explicações adicionais ou a seção 1 desse documento.

As seções desse documento contemplam os ambientes e/ou comandos mais relevantes nesse modelo, sempre seguidos de um exemplo de código e seu resultado.